% from ega0-16-fr.tex
\oldpage[0\textsubscript{IV-1}]{31}

\begin{proposition}[16.3.8]
\label{0.16.3.8-fr}
Soient $A$ un anneau semi-local noethérien, $X=\operatorname{Spec}(A)$ son
spectre, $\mathfrak{a}$ un idéal de $A$ distinct de $A$ tel que
$\operatorname{codim}(V(\mathfrak{a}),X)=r>0$. Il existe alors $r$ éléments
$x_1,\ldots,x_r$ de $\mathfrak{a}$, faisant partie d'un système de paramètres
de $A$, tels que
\[
  \operatorname{codim}(V(\mathfrak{a}),
  V(Ax_1+\cdots+Ax_r))=0.
\]
\end{proposition}

\begin{proof}
Soient $\mathfrak{p}_i$ $(1\leq i\leq m)$ les idéaux premiers minimaux de
l'anneau $A$, $\mathfrak{q}_j$ $(1\leq j\leq n)$ les idéaux premiers minimaux
parmi ceux qui contiennent $\mathfrak{a}$~; on a (14.2.1)
\[
  \operatorname{codim}(V(\mathfrak{a}),X)
  =\inf_j(\operatorname{codim}(V(\mathfrak{q}_j),X))
  =\inf_j(\dim(A_{\mathfrak{q}_j}))
\]
(16.1.3.2). L'hypothèse $r>0$ entraîne que $\mathfrak{a}$ n'est contenu dans
aucun des $\mathfrak{p}_i$~; par suite il existe $x\in\mathfrak{a}$ qui
n'appartient à aucun des $\mathfrak{p}_i$ (Bourbaki, \emph{Alg.~comm.},
chap.~II, \S~1, no~1, prop.~2). On a comme ci-dessus
\[
  \operatorname{codim}(V(\mathfrak{a}),V(Ax))
  =\inf_j(\dim((A/Ax)_{\mathfrak{q}_j/Ax}))
  =\inf_j(\dim(A_{\mathfrak{q}_j}/A_{\mathfrak{q}_j}x))~;
\]
mais comme $x/1$ n'est contenu dans aucun des idéaux premiers minimaux de
$A_{\mathfrak{q}_j}$, on a
$\dim(A_{\mathfrak{q}_j}/A_{\mathfrak{q}_j}x)
=\dim(A_{\mathfrak{q}_j})-1$ (16.3.4), d'où
$\operatorname{codim}(V(\mathfrak{a}),V(Ax))=r-1$. Pour la même raison, on a
$\dim(A/Ax)=\dim(A)-1$. Raisonnons maintenant par récurrence sur $r$~; si
$r>1$, il existe $r-1$ éléments $x'_i$ $(2\leq i\leq r)$ de $A'=A/Ax$,
faisant partie d'un système de paramètres de cet anneau, appartenant à
$\mathfrak{a}/Ax$ et tels que
\[
  \operatorname{codim}(V(\mathfrak{a}/Ax),
  V(A'x'_2+\cdots+A'x'_r))=0.
\]
Soit $(x'_i)_{r+1\leq i\leq s}$ un système d'éléments de $A'$ tel que
$(x'_i)_{2\leq i\leq s}$ soit un système de paramètres de $A'$~; pour tout
$i$ tel que $2\leq i\leq s$, soit $x_i$ un élément de la classe $x'_i$ dans
$A$, avec $x_i\in\mathfrak{a}$ pour $2\leq i\leq r$~; il est immédiat que
$A/(Ax+Ax_2+\cdots+Ax_s)$ est de longueur finie, et comme $\dim(A)=s$, on voit
que $x_1=x$ et les $x_i$ d'indices $i\geq2$ répondent aux conditions de
l'énoncé.
\end{proof}

\begin{proposition}[16.3.9]
\label{0.16.3.9-fr}
Soient $A$, $B$ deux anneaux locaux noethériens, $\mathfrak{m}$ l'idéal maximal
de $A$, $k$ son corps résiduel, $\varphi:A\to B$ un homomorphisme local.
\begin{enumerate}
  \item[\rm{(i)}] On a
  \begin{equation}
  \label{0.16.3.9.1-fr}
  \tag{16.3.9.1}
    \dim(B)\leq\dim(A)+\dim(B\otimes_A k).
  \end{equation}
  \item[\rm{(ii)}] Soient $M\neq0$ un $A$-module de type fini, $N\neq0$ un
  $B$-module de type fini. On a
  \begin{equation}
  \label{0.16.3.9.2-fr}
  \tag{16.3.9.2}
    \dim_B(M\otimes_A N)
    \leq\dim_A(M)+\dim_{B\otimes_A k}(N\otimes_A k).
  \end{equation}
\end{enumerate}
\end{proposition}

\begin{proof}
Soient $m$ la dimension de $A$, $(s_i)_{1\leq i\leq m}$ un système de
paramètres de $A$ (16.3.6), de sorte que si
$\mathfrak{a}=\sum_i s_iA$, $A/\mathfrak{a}$ est un $A$-module de longueur
finie. Alors l'anneau $A/\mathfrak{a}$ est artinien, donc l'idéal maximal
$\mathfrak{m}/\mathfrak{a}$ de cet anneau est nilpotent~; l'idéal
$\mathfrak{m}B/\mathfrak{a}B$ de $B/\mathfrak{a}B$, image de
$(\mathfrak{m}/\mathfrak{a})\otimes_A B$, est donc lui aussi nilpotent, et par
suite on a
$\dim(B/\mathfrak{a}B)=\dim(B/\mathfrak{m}B)$ (les spectres de ces deux
anneaux ayant même espace sous-jacent). D'autre part,
$B/\mathfrak{a}B=B/(s_1B+\cdots+s_mB)$, et les images des $s_i$ dans $B$
appartiennent à l'idéal maximal de $B$. Par suite (16.3.7), on a
\[
  \dim(B)\leq m+\dim(B\otimes_A k),
\]
ce qui n'est autre que (16.3.9.1).

Pour prouver (16.3.9.2), notons que si $\mathfrak{r}$ (resp.~$\mathfrak{s}$)
est l'annulateur de $M$ (resp.~$N$)
